\algrenewcommand\Return{\State \algorithmicreturn{} }%
\newcommand*\Let[2]{\State #1 $\gets$ #2} %atribuição

\usepackage{xparse}
\makeatletter
\NewDocumentCommand{\LeftComment}{s m}{%
  \Statex \IfBooleanF{#1}{\hspace*{\ALG@thistlm}}\(\triangleright\) #2}

\algdef{SE}[DOWHILE]{Do}{doWhile}{\algorithmicdo}[1]{\algorithmicwhile\ #1} %define do-while
\renewcommand{\algorithmicrequire}{\textbf{Requerimento:}}
\renewcommand{\algorithmicensure}{\textbf{Observação:}}

\renewcommand{\algorithmicreturn}{\textbf{retornar}}
\renewcommand{\algorithmicfunction}{\textbf{função}}
\renewcommand{\algorithmicprocedure}{\textbf{procedimento}}
\newcommand{\noProcedure}[2]{\Statex{\textbf{procedimento }#1(#2)}}
\newcommand{\noFunction}[2]{\Statex{\textbf{função }#1(#2)}}


%% Miscelânea: inicializadores e terminadores de blocos, comandos diversos

\renewcommand{\algorithmicdo}{\textbf{fa\c{c}a}}
\renewcommand{\algorithmicend}{\textbf{fim}}

%% Estruturas condicionais

\renewcommand{\algorithmicif}{\textbf{se}}
\renewcommand{\algorithmicthen}{\textbf{ent\~{a}o}}
\renewcommand{\algorithmicelse}{\textbf{sen\~{a}o}}

%% Estruturas de repetição

\renewcommand{\algorithmicwhile}{\textbf{enquanto}}

\renewcommand{\algorithmicfor}{\textbf{para}}
\renewcommand{\algorithmicforall}{\textbf{para todo}}

\renewcommand{\algorithmicloop}{\textbf{loop}}

\renewcommand{\algorithmicrepeat}{\textbf{repita}}
\renewcommand{\algorithmicuntil}{\textbf{at\'{e}}}

%Para cada (Foreach), adicionado por Juliana Souza
\algnewcommand\algorithmicforeach{\textbf{para cada}}
\algdef{S}[FOR]{ForEach}[1]{\algorithmicforeach\ #1\ \algorithmicdo}